\documentclass[14pt,a4paper]{extarticle}
\usepackage[T2A]{fontenc}
\usepackage[utf8]{inputenc}
\usepackage[english,russian]{babel}
\usepackage{lmodern}
\usepackage{geometry}
\usepackage{array}
\usepackage{tabularx}
\usepackage{longtable}
\usepackage{ltablex}
\usepackage{booktabs}
\usepackage{hyperref}
\usepackage{indentfirst}
\usepackage{microtype}
\usepackage{graphicx}
\usepackage{setspace}
\usepackage{listings}
\usepackage{graphicx}
\usepackage{caption}
\usepackage{float}
\usepackage{tocloft}
\usepackage{titlesec}
\usepackage{newtxtext,newtxmath}

\titleformat{\section}        % форматируем \section
  {\normalfont\bfseries}      % стиль: обычный + жирный
  {\thesection.}              % номер секции с точкой
  {2mm}                       % расстояние между номером и текстом
  {}                           % перед заголовком можно добавить что-то
\titlespacing*{\section}{1.5cm}{*3}{*1}  % {отступ слева}{отступ сверху}{отступ снизу}

\captionsetup{justification=centering}
\renewcommand{\ttdefault}{cmtt}
\keepXColumns
\setlength{\parindent}{3em}
\newgeometry{left=2.5cm,right=2.5cm,top=3cm,bottom=2.5cm}

\renewcommand{\contentsname}{\hspace*{4em}{\smallСодержание}}
\setlength{\cftbeforesecskip}{0.5em}
\renewcommand{\cftsecleader}{\cftdotfill{\cftdotsep}}
\renewcommand{\cftsecaftersnum}{.\hspace{2mm}}

\title{Отчёт по "Организация работы ИТ-службы"}
\author{Выполнил студент группы М-РИТ-25-1, Станиславчук Сергей}

\begin{document}

\setstretch{1.5}
\noindent
\begin{minipage}[t]{0.01\textwidth}
    \includegraphics{~/programmer/++Programmer/5_1/LSTU.jpg}
\end{minipage}%
\begin{minipage}[c]{1.1\linewidth}
    \centering
    {\small
    МИНИСТЕРСТВО НАУКИ И ВЫСШЕГО ОБРАЗОВАНИЯ РФ \\[1em]
    ФЕДЕРАЛЬНОЕ ГОСУДАРСТВЕННОЕ БЮДЖЕТНОЕ \\
    ОБРАЗОВАТЕЛЬНОЕ УЧРЕЖДЕНИЕ ВЫСШЕГО ОБРАЗОВАНИЯ \\
    «ЛИПЕЦКИЙ ГОСУДАРСТВЕННЫЙ ТЕХНИЧЕСКИЙ УНИВЕРСИТЕТ»
    }
\end{minipage}%

\vspace{3em}

\noindent
\hspace{1em}%
\begin{minipage}[t]{0.8\textwidth}
\textbf{Институт} \hspace{2em} компьютерных наук \\
\textbf{Кафедра} \hspace{2em} автоматизированных систем управления
\end{minipage}


\vspace{6em}

\begin{center}
\textbf{Практическая работа №1} \\
по дисциплине \\ <<Организация работы ИТ-службы>>
\end{center}

\vspace{6em}

\noindent
\begin{tabular*}{\textwidth}{@{\extracolsep{\fill}} l c r }
    Студент М-РИТ-25-1 & \makebox[4cm]{\hrulefill} & Станиславчук С.М. \\[2em]
Руководитель & \\
К.т.н., доцент & \makebox[4cm]{\hrulefill} & Ведищев В.В. \\
\end{tabular*}

\vfill
\begin{center}
    Липецк, 2026
\end{center}

\thispagestyle{empty}

\newpage
\tableofcontents
\newpage

\section{Выбор сервиса}

В качестве рассматриваемого IT-сервиса выбрана услуга предоставления виртуальных выделенных серверов (VDS/VPS). Данный сервис предполагает предоставление заказчику виртуальной вычислительной среды с выделенными ресурсами процессора, оперативной памяти, дискового пространства и сетевого подключения.

Сервис предоставляет клиентам виртуальные серверы на собственной облачной платформе. При этом физическая инфраструктура может размещаться в дата-центрах, а сетевое подключение обеспечивается с использованием внешних операторов связи. Таким образом, при предоставлении услуги участвуют как непосредственно поставщик сервиса, так и внешние поставщики инфраструктурных компонентов.

Основными заказчиками услуги являются организации и частные лица, которым требуется вычислительный ресурс для размещения собственных приложений, веб-сайтов, баз данных и других информационных систем без необходимости приобретения и самостоятельного обслуживания физического сервера.

Пользователями сервиса могут выступать системные администраторы, разработчики, DevOps-инженеры и другие специалисты, осуществляющие эксплуатацию размещенных на сервере информационных систем.

Таким образом, в рамках данной работы под рассматриваемым IT-сервисом понимается предоставление и эксплуатация виртуального сервера для заказчика.

\section{Стратегия сервиса}

Стратегия сервиса определяет назначение услуги, потребности заказчиков и создаваемую для них ценность. Основная ценность VDS заключается в предоставлении заказчику вычислительных ресурсов без необходимости приобретения физического сервера и самостоятельного обеспечения всей инфраструктуры, необходимой для его работы.

Использование VDS позволяет заказчику получить изолированную виртуальную вычислительную среду с заданным количеством процессорных ресурсов, оперативной памяти, дискового пространства и сетевым подключением. При этом обслуживание физической инфраструктуры, на которой работает виртуальный сервер, остается зоной ответственности поставщика.

Основными заинтересованными сторонами сервиса являются заказчики и пользователи VDS, поставщик услуги, сотрудники технической поддержки и эксплуатации, а также внешние поставщики инфраструктуры. Интересы этих участников могут различаться. Заказчику необходимы доступность, производительность и предсказуемая стоимость услуги, тогда как поставщик должен одновременно обеспечивать качество сервиса и рациональное использование собственной инфраструктуры.

Для оценки и дальнейшего управления услугой используются три уровня управления. Стратегический уровень определяет направление развития сервиса и его соответствие потребностям рынка. Тактический уровень определяет способы организации предоставления услуги, необходимые ресурсы и параметры качества. Операционный уровень обеспечивает непосредственное функционирование инфраструктуры и обработку возникающих событий. Более подробно данные уровни управления рассматриваются в последующих практических работах.

Стратегическая цель сервиса заключается в предоставлении заказчику доступной и управляемой вычислительной инфраструктуры, позволяющей ему сосредоточиться на эксплуатации собственных информационных систем, не приобретая и не обслуживая физические серверы.

\section{Проектирование сервиса}

Проектирование VDS-сервиса включает определение его функциональности, архитектуры, параметров качества, необходимых ресурсов и процедур предоставления услуги.

Основной функциональностью сервиса является создание для заказчика виртуальной машины с заданными вычислительными и сетевыми характеристиками. В состав услуги входят процессорные ресурсы, оперативная память, дисковое пространство и сетевое подключение.

Архитектура сервиса включает несколько основных компонентов. На физическом уровне используются серверы, системы хранения данных и сетевое оборудование. На программном уровне используется платформа виртуализации и программное обеспечение управления инфраструктурой. На созданной виртуальной машине размещается операционная система и программное обеспечение заказчика.

Упрощенно взаимодействие компонентов можно представить следующим образом:

\begin{center}
Физическая инфраструктура $\rightarrow$ платформа виртуализации $\rightarrow$ виртуальный сервер $\rightarrow$ операционная система $\rightarrow$ приложения заказчика
\end{center}

Основными параметрами IT-сервиса являются функциональность, время обслуживания, доступность, надежность, производительность, конфиденциальность, масштаб и затраты.

Время обслуживания для VDS должно соответствовать характеру услуги. Поскольку виртуальный сервер используется для размещения постоянно работающих информационных систем, целевым режимом является круглосуточное предоставление услуги, то есть 24/7.

Доступность характеризует долю согласованного времени, в течение которого VDS доступен заказчику. Для оценки данного параметра может использоваться целевое значение SLA, например 99,9% доступности. Фактическое значение должно определяться по результатам мониторинга.

Производительность определяется характеристиками выделенного виртуального сервера. К основным показателям относятся количество виртуальных процессорных ядер, объем оперативной памяти, производительность дисковой подсистемы и пропускная способность сетевого подключения.

Надежность сервиса обеспечивается резервированием инфраструктуры, мониторингом состояния физических серверов и систем хранения, а также процедурами восстановления после отказов. Конкретные методы обеспечения надежности зависят от архитектуры инфраструктуры поставщика.

Конфиденциальность обеспечивается изоляцией виртуальных серверов друг от друга, механизмами аутентификации и авторизации, сетевыми средствами защиты и соблюдением требований информационной безопасности. При этом заказчик самостоятельно отвечает за безопасность программного обеспечения и данных, размещенных внутри принадлежащего ему виртуального сервера, в пределах распределения ответственности между заказчиком и поставщиком.

Для предоставления сервиса необходимы следующие ресурсы:

\begin{itemize}
\item физические серверы;
\item оперативная память и процессорные ресурсы;
\item системы хранения данных;
\item сетевое оборудование и каналы связи;
\item программная платформа виртуализации;
\item системы мониторинга;
\item системы резервного копирования;
\item средства информационной безопасности;
\item персонал, осуществляющий эксплуатацию и техническую поддержку.
\end{itemize}

Стоимость услуги должна учитывать затраты на приобретение и эксплуатацию оборудования, размещение инфраструктуры, электроэнергию, сетевые услуги, программное обеспечение и оплату труда сотрудников.

При проектировании сервиса также необходимо учитывать требования к допустимому использованию инфраструктуры. Поставщик устанавливает правила использования услуги, запрещающие применение серверов для противоправной деятельности, атак на другие системы, распространения вредоносного программного обеспечения и других действий, нарушающих условия предоставления услуги или законодательство.

\section{Внедрение (преобразование) сервиса}

Внедрение VDS-сервиса представляет собой комплекс мероприятий, в результате которых спроектированная услуга становится доступной для пользователей.

На первоначальном этапе определяется необходимая производительность и масштаб инфраструктуры. После этого выполняются подготовка или закупка физических серверов, систем хранения и сетевого оборудования, настройка сетевой инфраструктуры и установка программной платформы виртуализации.

Следующим этапом является настройка средств управления виртуальными серверами. Они должны обеспечивать создание, изменение, остановку и удаление виртуальных машин, распределение вычислительных ресурсов и контроль состояния инфраструктуры.

После подготовки основной инфраструктуры выполняется настройка систем мониторинга и резервного копирования. Мониторинг позволяет обнаруживать отказы и снижение качества сервиса, а резервное копирование обеспечивает возможность восстановления данных в предусмотренных сценариях отказа.

До начала промышленной эксплуатации проводится тестирование. Проверяется корректность создания виртуальных серверов, доступность сетевых ресурсов, производительность, изоляция виртуальных машин, корректность резервного копирования и восстановление после отказов.

После завершения тестирования сервис передается в эксплуатацию. Для него определяются правила предоставления услуги, порядок обработки обращений пользователей, процедуры реагирования на инциденты и порядок внесения изменений.

Таким образом, результатом этапа преобразования является готовый к предоставлению сервис с подготовленной технической инфраструктурой, регламентами эксплуатации и определенными параметрами качества.

\section{Эксплуатация сервиса}

Эксплуатация является наиболее продолжительным этапом жизненного цикла VDS-сервиса. На данном этапе поставщик должен обеспечивать непрерывное функционирование услуги в соответствии с установленными требованиями и согласованными уровнями обслуживания.

Основным инструментом эксплуатации является мониторинг. Контролируются доступность виртуальных серверов, загрузка процессора, использование оперативной памяти, состояние дисковой подсистемы, сетевой трафик и состояние физических компонентов инфраструктуры.

В процессе эксплуатации возникают обращения пользователей. Например, заказчик может сообщить о недоступности виртуального сервера, снижении производительности или проблемах с сетевым подключением. Такое обращение регистрируется и обрабатывается в соответствии с установленной процедурой.

Отдельную группу составляют запросы на обслуживание. К ним можно отнести изменение конфигурации виртуального сервера, увеличение объема оперативной памяти или дискового пространства, изменение сетевых параметров и другие действия, предусмотренные условиями предоставления услуги.

Неисправность физического сервера может привести к нарушению работы нескольких виртуальных серверов. В этом случае задача поставщика заключается в локализации причины неисправности, восстановлении работоспособности сервиса и информировании затронутых заказчиков.

Изменения инфраструктуры также являются частью эксплуатации. К ним относятся обновление программного обеспечения виртуализации, замена неисправного оборудования, модернизация систем хранения и изменение сетевой инфраструктуры. Такие изменения должны выполняться контролируемым образом, чтобы минимизировать влияние на работающие виртуальные серверы.

Важной частью эксплуатации является обеспечение информационной безопасности. Поставщик должен реагировать на нарушения правил использования сервиса. Если клиент использует VDS для запрещенной деятельности, могут применяться предусмотренные правилами провайдера меры, включая ограничение или приостановку обслуживания. При наличии соответствующих оснований информация может быть предоставлена компетентным органам в установленном порядке.

\section{Постоянное улучшение сервиса}

VDS-сервис, как и любая информационная система, функционирует в условиях постоянного изменения требований пользователей, технологий и внешней среды. Поэтому после ввода услуги в эксплуатацию необходимо регулярно оценивать ее качество и определять направления дальнейшего развития.

Основой для улучшения являются показатели эксплуатации. Анализируются доступность сервиса, количество и продолжительность инцидентов, время восстановления, производительность инфраструктуры, количество обращений пользователей и затраты на предоставление услуги.

Отдельное значение имеет анализ повторяющихся инцидентов. Если одна и та же неисправность возникает регулярно, необходимо определить ее первопричину и принять меры, предотвращающие повторение проблемы. Например, систематические отказы определенного оборудования могут привести к его замене или изменению архитектуры инфраструктуры.

Развитие технологий также является основанием для модернизации. Появление более производительных процессоров, накопителей и сетевого оборудования позволяет увеличивать производительность сервиса и расширять доступный объем ресурсов.

Изменение потребностей заказчиков может привести к изменению состава и параметров тарифов. Например, при росте спроса на вычислительные ресурсы могут быть введены новые конфигурации виртуальных серверов с увеличенным объемом оперативной памяти, количеством виртуальных процессоров или дискового пространства.

Постоянное улучшение может затрагивать не только техническую инфраструктуру, но и организационные процессы. Могут совершенствоваться процедуры обработки обращений, мониторинга, резервного копирования, реагирования на инциденты и информационной безопасности.

Таким образом, постоянное улучшение направлено на поддержание соответствия VDS-сервиса потребностям заказчиков и развитие инфраструктуры поставщика. Результаты анализа эксплуатации используются для определения изменений, которые должны быть реализованы в следующих циклах развития сервиса.

\end{document}
